\documentclass[10pt,a4paper]{amsart}
\usepackage[margin=19mm]{geometry}
\usepackage{amsmath,amssymb,mathtools}
\usepackage[scr=rsfs,cal=euler]{mathalfa}
\usepackage{xfrac}
\usepackage[unicode,hidelinks,bookmarks=false]{hyperref}
\newcommand{\ie}{i.e.\ }
\newcommand{\cf}{cf.\ }
\newcommand{\resp}{respectively\ }
\newcommand{\Subsection}{\subsection}
\providecommand{\nobreakdash}{\nobreak\hbox{-}}
\setlength{\emergencystretch}{2em}
\multlinegap0pt
\usepackage{amssymb}
\newtheorem{lem}{Lemma}
\newcommand{\nn}{\nonumber}
\title{On determinantal formulas for hermitian random matrices}
\author{Di YANG, Jiayi ZHAO, Jian ZHOU}
\date{Source: arXiv:2606.11958v1 (CC BY 4.0)}
\begin{document}
\maketitle

\noindent\emph{Source and licence.} On determinantal formulas for hermitian random matrices. Version arXiv:2606.11958v1 (CC BY 4.0), distributed by arXiv under the Creative Commons Attribution 4.0 International licence, http://creativecommons.org/licenses/by/4.0/, as recorded in the arXiv licence metadata for this exact version and shown on the abstract page https://arxiv.org/abs/2606.11958v1. Full licence notice: \texttt{licenses/arxiv-2606.11958-CC-BY-4.0.txt}.\par\smallskip

\noindent\textit{Editorial scope.} Counted unit: the article's lem lk (source0.tex lines 547--553). The article's complete original proof of it is delivered with it (source0.tex lines 554--665, one \texttt{proof} environment of 112 lines). No mathematics is written, repaired or re-derived: the statement and the proof are the article's own lines, in the article's own notation, and no retained line is substituted or re-flowed.  Retained uncounted as the source-local context the proof needs: lines 528--539.  Nothing was omitted from any retained span, so the package carries no omission record.  No retained line contains a citation, so the wrapper carries no bibliography and no bibliography entry.  No retained line contains a figure environment, an image-inclusion command or drawing code, so no image is delivered and the package declares no asset dependency.  Editorial formatting only, in addition: an AMS \texttt{amsart} preamble that restates the article's own declarations and macros used by the retained lines, namely the environments \texttt{lem} on separate counters, together with the standard packages those lines need.  The retained environments print in this document as Lemma 1 in order of appearance, whereas the article prints the same items with the numbering of its own full document.  The complete native source of the article as distributed in the arXiv e-print for this exact version is bundled unchanged as \texttt{sources/202588/source0.tex}.
	\begin{equation} \label{table}
		\begin{array}{c|cccccc}
			\text{row}\backslash\text{column} & 0 & 1 & 2 & \dots & k-1 & k \\
			\hline
			0 & \ell_0 & \ell_1 & \ell_2 & \dots & \ell_{k-1} & \ell_k \\
			1 & & \ell_{1,(1)} & \ell_{2,(1)} & \dots & \ell_{k-1,(1)} & \ell_{k,(1)} \\
			2 & & & \ell_{2,(2)} & \dots & \ell_{k-1,(2)} & \ell_{k,(2)} \\
			\vdots & & & & \ddots & \vdots & \vdots \\
			k-1 & & & & & \ell_{k-1,(k-1)} & \ell_{k,(k-1)} \\
			k & & & & & & \ell_{k,(k)}
		\end{array}
	\end{equation} 

\begin{lem}
	For $k\ge2$,
	\begin{align}
		\ell_k(n;\lambda_1,\cdots,\lambda_k)=&(-1)^k\prod_{i=1}^k\mathcal J_n(\lambda_{i+1},\lambda_i)-\prod_{i=1}^k\frac{1}{(\lambda_i-\lambda_{i+1})}\nn\\
		&-\sum_{t=1}^{k-1}\sum_{1\le a_1<\cdots<a_t\le k}R_{a_1,\cdots,a_t}(\lambda_1,\cdots,\lambda_k)\ell_t(n;\lambda_{a_1},\cdots,\lambda_{a_t}). \label{lk}
	\end{align}
\end{lem}

\begin{proof}
	We first outline the strategy of the proof. In the context of Table~\eqref{table}, equation~\eqref{lk} can be interpreted as expressing the entry at $(0,k)$ in terms of the entries at $(0,m)$, $0\le m\le k-1$ and the diagonal entry at $(k,k)$. To achieve this, we initially express the entry at $(0,k)$ in terms of the entries at the $(k-1)$th column and at the diagonal term $(k,k)$. Subsequently, in each step, we iteratively reduce the dependence on the $m$-th column to the $(m-1)$-th column and the boundary entry at $(0,m)$. This procedure ultimately yields~\eqref{lk}. The process is illustrated as follows:
	\[
	\begin{array}{c|c}
		 & k \\ \hline
		 0 & *
	\end{array}
	\Rightarrow
	\begin{array}{c}
		\begin{array}{c|cc}
			 & k-1 & k \\ \hline
			0 & * & * \\
			1 & * & 0 \\
			\vdots & \vdots & \vdots \\
			k-1 & * & 0 \\
			k & & * 
		\end{array}
	\end{array}
	\Rightarrow
	\begin{array}{c}
		\begin{array}{c|ccc}
			& k-2 & k-1 & k \\ \hline
			0 & * & * & * \\
			1 & * & 0 & 0 \\
			\vdots & \vdots & \vdots & \vdots \\
			k-2 & * & 0 & 0 \\
			k-1 & & 0 & 0 \\
			k & &  & * 
		\end{array}
	\end{array}
	\Rightarrow \dots \Rightarrow
	\begin{array}{c}
		\begin{array}{c|ccccc}
			& 0 & 1 & \dots & k-1 & k \\ \hline
			0 & * & * & \dots & * & * \\
			1 & & 0 & \dots & 0 & 0 \\
			\vdots & &  & \ddots & \vdots & \vdots \\
			k-1 & &  &  & 0 & 0 \\ 
			k & & &  & & * 
		\end{array}
	\end{array}
	\]
	
	More precisely, for $m\le k-1$ we prove the following formula, which expresses $\ell_k$ in terms of the entries at $(0,r)$ for $k-m\le r\le k-1$, the entries at $(r,k-m)$ for $1\le r\le k-m$, and the diagonal term $(k,k)$: (We still omit the parameter $n$.)
	\begin{align}
		&\ell_k(\lambda_1,\cdots,\lambda_k)\nn \\
		=& \ell_{k,(k)}-\sum_{r=1}^{m-1}\sum_{1\le a_1<\cdots<a_r\le k}\widetilde R_{a_1,\cdots,a_r}\ell_{k-r}^{\widehat{a_1,\cdots,a_r}} \nn\\
		& -\sum_{r=0}^{m-1}\sum_{r+1<a_{r+1}<\cdots<a_{m-1}\le k-1}\widetilde R_{1,\cdots,r,a_{r+1},\cdots,a_{m-1}}\frac{\ell_{k-m}^{\widehat{1,\cdots,r,a_{r+1},\cdots,a_{m-1},k}}-\ell_{k-m}^{\widehat{1,\cdots,r,r+1,a_{r+1},\cdots,a_{m-1}}}}{\lambda_k-\lambda_1} \nn\\
		& -\sum_{t=1}^{k-m-1}\sum_{r=0}^{m-1}\sum_{r+1<a_{r+1}<\cdots<a_m=k-t} \frac{\lambda_k-\lambda_{r+1}}{\lambda_k-\lambda_1}\widetilde R_{1,\cdots,r,a_{r+1},\cdots,a_m}\ell_{k-m,(t)}^{\widehat{1,\cdots,r,a_{r+1},\cdots,a_m}}\nn\\
		&-\frac{\lambda_k-\lambda_{m+1}}{\lambda_k-\lambda_1}\widetilde R_{1,\cdots,m}\ell_{k-m,(k-m)}^{\widehat{1,\cdots,m}}, \label{lind}
	\end{align}
	where for $a_{1}<\cdots<a_{m}$ we denote
	\begin{align}
		\widetilde R_{a_1,\cdots,a_m}:=& \prod_{s=1}^{m}\frac{1}{\lambda_{a_s}-\lambda_{a_s+1}}, \\
		\ell_{k-m}^{\widehat{a_1,\cdots,a_m}} :=& \ell_{k-m}(\lambda_1,\cdots,\lambda_{a_1-1},\lambda_{a_1+1},\cdots,\lambda_{a_m-1},\lambda_{a_m+1},\cdots,\lambda_k), \\
		\ell_{k-m,(t)}^{\widehat{a_1,\cdots,a_m}} :=& \ell_{k-m,(t)}(\lambda_1,\cdots,\lambda_{a_1-1},\lambda_{a_1+1},\cdots,\lambda_{a_m-1},\lambda_{a_m+1},\cdots,\lambda_k).
	\end{align}
	
	We prove~\eqref{lind} by induction on $m$. First, for $m=1$,
	\begin{align}
		\ell_{k}(\lambda_1,\cdots,\lambda_k) &= \ell_{k,(1)}-\frac{1}{\lambda_k-\lambda_1}(\ell_{k-1}^{\widehat k}-\ell_{k-1}^{\widehat1}) \nn\\
		&= \ell_{k,(2)}-\frac{1}{\lambda_k-\lambda_1}(\ell_{k-1}^{\widehat k}-\ell_{k-1}^{\widehat1})-\frac{1}{\lambda_{k-1}-\lambda_k}\ell_{k-1,(1)}^{\widehat{k-1}} \nn\\
		&= \ell_{k,(k-1)}-\frac{1}{\lambda_k-\lambda_1}(\ell_{k-1}^{\widehat k}-\ell_{k-1}^{\widehat1})-\sum_{t=1}^{k-2}\frac{1}{\lambda_{k-t}-\lambda_{k-t+1}}\ell_{k-1,(t)}^{\widehat{k-t}} \nn\\
		&= \ell_{k,(k)}-\frac{1}{\lambda_k-\lambda_1}(\ell_{k-1}^{\widehat k}-\ell_{k-1}^{\widehat1})-\sum_{t=1}^{k-2}\widetilde R_{k-t}\ell_{k-1,(t)}^{\widehat{k-t}} -\frac{\lambda_k-\lambda_2}{\lambda_k-\lambda_1}\widetilde R_1\ell_{k-1,(k-1)}^{\widehat1}.
	\end{align}
	
	Now assume that~\eqref{lind} holds for $m$. For $m+1$, we transition from the $(k-m)$-th column to the $(k-m-1)$-th column. Consequently, we need to compute the terms at the entries $(t,k-m-1)$ for $0\le t\le k-m-1$, as well as the boundary term $(0,k-m)$. The term at $(k-m-1,k-m-1)$ is derived from the term at $(k-m,k-m)$ and is given by
	\begin{align}
		&-\left(\frac{1}{\lambda_{m+1}-\lambda_{m+2}}+\frac{1}{\lambda_k-\lambda_{m+1}}\right)\frac{\lambda_k-\lambda_{m+1}}{\lambda_k-\lambda_1}\widetilde R_{1,\cdots,m}\ell_{k-m-1,(k-m-1)}^{\widehat{1,\cdots,m,m+1}}\nn\\
		=&-\frac{\lambda_k-\lambda_{m+2}}{\lambda_k-\lambda_1}\widetilde R_{1,\cdots,m+1}\ell_{k-m-1,(k-m-1)}^{\widehat{1,\cdots,m+1}}. \label{part1}
	\end{align}
	The term at $(t,k-m-1)$ for $1\le t\le k-m-2$ is
	\begin{align}
		&-\frac{1}{\lambda_{k-t}-\lambda_{k-t+1}}\sum_{r=0}^{m}\sum_{r+1<a_{r+1}<\cdots<a_m\le k-t-1}\frac{\lambda_k-\lambda_{r+1}}{\lambda_k-\lambda_1}\widetilde R_{1,\cdots,r,a_{r+1},\cdots,a_m}\ell_{k-m-1,(t)}^{\widehat{1,\cdots,r,a_{r+1},\cdots,a_m,k-t}}\nn\\
		=&-\sum_{r=0}^{m+1}\sum_{r+1<a_{r+1}<\cdots<a_{m+1}=k-t}\frac{\lambda_k-\lambda_{r+1}}{\lambda_k-\lambda_1}\widetilde R_{1,\cdots,r,a_{r+1},\cdots,a_{m+1}}\ell_{k-m-1,(t)}^{\widehat{1,\cdots,r,a_{r+1},\cdots,a_{m+1}}}. \label{part2}
	\end{align}
	Similarly, the term at $(0,k-m-1)$ reads
	\begin{align}
		&-\sum_{r=0}^{m}\sum_{r+1<a_{r+1}<\cdots<a_m\le k-1}\frac{\lambda_k-\lambda_{r+1}}{\lambda_k-\lambda_1}\widetilde R_{1,\cdots,r,a_{r+1},\cdots,a_m}\frac{\ell_{k-m-1}^{\widehat{1,\cdots,r,a_{r+1},\cdots,a_m,k}}-\ell_{k-m-1}^{\widehat{1,\cdots,r+1,a_{r+1},\cdots,a_m}}}{\lambda_k-\lambda_{r+1}}\nn\\
		=&-\sum_{r=0}^{m}\sum_{r+1<a_{r+1}<\cdots<a_m\le k-1}\widetilde R_{1,\cdots,r,a_{r+1},\cdots,a_m}\frac{\ell_{k-m-1}^{\widehat{1,\cdots,r,a_{r+1},\cdots,a_m,k}}-\ell_{k-m-1}^{\widehat{1,\cdots,r+1,a_{r+1},\cdots,a_m}}}{\lambda_k-\lambda_1}. \label{part3}
	\end{align}
	Finally, the term at $(0,k-m)$ changes to
	\begin{align}
		& -\sum_{r=0}^{m-1}\sum_{r+1<a_{r+1}<\cdots<a_{m-1}\le k-1}\widetilde R_{1,\cdots,r,a_{r+1},\cdots,a_{m-1}}\frac{\ell_{k-m}^{\widehat{1,\cdots,r,a_{r+1},\cdots,a_{m-1},k}}-\ell_{k-m}^{\widehat{1,\cdots,r+1,a_{r+1},\cdots,a_{m-1}}}}{\lambda_k-\lambda_1} \nn\\
		&\qquad -\sum_{t=1}^{k-m}\sum_{r=0}^{m}\sum_{r+1<a_{r+1}<\cdots<a_m=k-t} \frac{\lambda_k-\lambda_{r+1}}{\lambda_k-\lambda_1}\widetilde R_{1,\cdots,r,a_{r+1},\cdots,a_m}\ell_{k-m}^{\widehat{1,\cdots,r,a_{r+1},\cdots,a_m}}\nn\\
		&\qquad -\frac{\lambda_k-\lambda_{m+1}}{\lambda_k-\lambda_1}\widetilde R_{1,\cdots,m}\ell_{k-m}^{\widehat{1,\cdots,m}}.\label{0k-m}
	\end{align}
	In~\eqref{0k-m} we compute the rational function before $\ell_{k-m}^{\widehat{1,\cdots,r,a_{r+1},\cdots,a_{m}}}$ where $a_{r+1}>r+1$ and $a_{m}<k$. It is
	\begin{align}
		&-\sum_{t=0}^{r-1}\frac{1}{\lambda_1-\lambda_k}\widetilde R_{1,\cdots,t,t+2,\cdots,r,a_{r+1},\cdots,a_{m-1}}-\frac{\lambda_k-\lambda_{r+1}}{\lambda_k-\lambda_1}\widetilde R_{1,\cdots,r,a_{r+1},\cdots,a_m} \nn\\
		=&-\frac{1}{\lambda_1-\lambda_k}\left(\sum_{t=0}^{r-1}(\lambda_{t+1}-\lambda_{t+2})-(\lambda_k-\lambda_{r+1})\right)\widetilde R_{1,\cdots,r,a_{r+1},\cdots,a_{m}} \nn\\
		=&-\widetilde R_{1,\cdots,r,a_{r+1},\cdots,a_m}.
	\end{align}
	Similarly, the rational function before $\ell_{k-m}^{\widehat{1,\cdots,m}}$ is $-\widetilde R_{1,\cdots,m}$. And lastly, the rational function before $\ell_{k-m}^{\widehat{1,\cdots,r,a_{r+1},\cdots,a_{m-1},k}}$ where $a_{r+1}>r+1$ is
	\begin{align}
	-\frac{1}{\lambda_k-\lambda_1}\widetilde R_{1,\cdots,r,a_{r+1},\cdots,a_{m-1}} =-\widetilde R_{1,\cdots,r,a_{r+1},\cdots,a_{m-1},k}.
	\end{align}
	Hence~\eqref{0k-m} reads
	\begin{equation}
		-\sum_{1\le a_1<\cdots<a_m\le k}\widetilde R_{a_1,\cdots,a_m}\ell_{k-m}^{\widehat{a_1,\cdots,a_m}}. \label{part4}
	\end{equation}
	Thus by~\eqref{part1},~\eqref{part2},~\eqref{part3} and~\eqref{part4}, we obtain~\eqref{lind} for $m+1$.
	
	Now following~\eqref{lind}, for $m=k-1$ we have
	\begin{align}
		&\ell_k(\lambda_1,\cdots,\lambda_k)\nn\\
		=& \ell_{k,(k)}-\sum_{r=1}^{k-2}\sum_{1\le a_1<\cdots<a_r\le k}\widetilde R_{a_1,\cdots,a_r}\ell_{k-r}^{\widehat{a_1,\cdots,a_r}}-\sum_{r=0}^{k-2}\widetilde R_{1,\cdots,r,r+2,\cdots,k-1}\frac{\ell_1(\lambda_{r+1})-\ell_1(\lambda_k)}{\lambda_k-\lambda_1} \nn\\
		=& \ell_{k,(k)}-\sum_{r=1}^{k-1}\sum_{1\le a_1<\cdots<a_r\le k}\widetilde R_{a_1,\cdots,a_r}\ell_{k-r}^{\widehat{a_1,\cdots,a_r}} \nn\\
		=&(-1)^k\prod_{i=1}^k\mathcal J_n(\lambda_{i+1},\lambda_i)-\frac{1}{\prod_{i=1}^k(\lambda_i-\lambda_{i+1})}-\sum_{t=1}^{k-1}\sum_{1\le a_1<\cdots<a_t\le k}R_{a_1,\cdots,a_t}(\lambda_1,\cdots,\lambda_k)\ell_t(\lambda_{a_1},\cdots,\lambda_{a_t}).
	\end{align}
	Thus the lemma is proved.
\end{proof}

\end{document}
